% TOP_PROBLEM: {"id": 2302089, "title": "Research Problems in Function Theory — Problem 2.89", "category_id": 8, "category": {"id": 8, "name": "analysis", "display_name": "Analysis", "description": "Limits, continuity, calculus, and function theory.", "slug": "analysis", "order_index": 8, "created_at": "2026-07-31T15:26:25.671Z"}, "status": "open"}
% FIRST_POSTED: null
% ENTRY_KIND: statement_only
% Imported from: batch12.tex; UnsolvedMath ID 2302089
% Compile twice: lualatex 2302089.tex
\documentclass[11pt]{article}
\usepackage{fontspec}
\setmainfont{Latin Modern Roman}
\usepackage{amsmath,amssymb,mathtools,unicode-math}
\setmathfont{Latin Modern Math}
\usepackage{xurl,hyperref}
\hypersetup{colorlinks=true,urlcolor=blue,pdftitle={UnsolvedMath Problem 2302089}}
\newcommand{\sref}[2]{\hyperlink{source:#1:#2}{[S#2]}}
\newcommand{\cataloguescope}[1]{\paragraph{Mathematical question.}#1}
\begin{document}
% UnsolvedMath ID 2302089; source catalogue code AMR-022-2089
\setcounter{section}{2302088}
\section{Research Problems in Function Theory — Problem 2.89}\label{2302089}
{\small\sffamily UnsolvedMath ID 2302089 · Analysis}
\subsection{Definitions and mathematical statement}

\paragraph{Shared notation.}$\mathbb N=\{1,2,\ldots\}$, and $\mathbb Z,\mathbb Q,\mathbb R,\mathbb C$ denote the integers, rationals, reals and complex numbers. Any other conventions are specified locally.


Let $\mathcal R_d$ be the space of rational maps of degree $d\geq2$, with its usual coefficient topology (pairs of degree-at-most-$d$ homogeneous polynomials, modulo a common nonzero scalar, with nonzero resultant). An invariant Herman ring is a doubly connected Fatou component on which the map is conformally conjugate to $w\mapsto e^{2\pi i\alpha}w$ on an annulus $1<|w|<R$. A real irrational $\alpha$ is Diophantine if some $C>0$ and $\tau>0$ satisfy $|\alpha-p/q|\geq Cq^{-(2+\tau)}$ for all integers $p$ and $q\geq1$.
Fix $f_0\in\mathcal R_d$ and an invariant Herman ring $A_0$ of Diophantine rotation number $\alpha$. Let $\mathcal H_d(\alpha;f_0,A_0)$ consist of endpoints $f_1$ of continuous paths $f_t\in\mathcal R_d$, $0\leq t\leq1$, along which there are invariant Herman rings $A_t$ of rotation number $\alpha$, starting with $A_0$, that vary continuously in the Carath\'eodory kernel sense (compact subsets of a limiting domain eventually belong to the corresponding domains, with the component determined by a moving interior base point).
\paragraph{Statement recorded in the supplied source.}
Is $\mathcal H_d(\alpha;f_0,A_0)$ locally closed in $\mathcal R_d$? Is its boundary relative to its closure in $\mathcal R_d$, namely $\partial_{\overline{\mathcal H_d(\alpha;f_0,A_0)}}\mathcal H_d(\alpha;f_0,A_0)$, a topological manifold? \sref{2302089}{2}
\subsection{Short English statement}
For a fixed Diophantine Herman-ring rotation, is the path-connected parameter locus locally closed, and is its boundary within its closure a manifold?
\subsection{Sources}
\begin{description}
\item[\hypertarget{source:2302089:1}{S1}] ULAM.AI and the UnsolvedMath Contributors, UnsolvedMath: A Curated Collection of Open Mathematics Problems, record 2302089 (AMR-022-2089). CC BY 4.0. \url{https://huggingface.co/datasets/ulamai/UnsolvedMath}.
\item[\hypertarget{source:2302089:2}{S2}] W. K. Hayman and E. F. Lingham, \emph{Research Problems in Function Theory}, arXiv:1809.07200v2 (2018), Problem 2.89. \url{https://arxiv.org/abs/1809.07200}.
\end{description}
\label{end:2302089}
\end{document}
